\chapter{Restriction Reference}
\label{ch:restrictions}

A one-page summary of \cref{sec:l0-restrictions}, for use while writing or
reviewing DSTM code. The statements here are abbreviations; the normative
\glspl{version} are \crefrange{res:monotone}{res:terminate}.

Throughout, $p$ is the issuing \gls{rank}, $q$ the target \gls{rank}, $t$ the \gls{origin-time},
$\delta$ the \gls{delta} and $\eff = t + \delta$ the \gls{effective-time}.

\begin{table}[htbp]
\centering
\small
\renewcommand{\arraystretch}{1.35}
\begin{tabularx}{\textwidth}{@{}l >{\raggedright\arraybackslash}p{4.3cm}
                                 >{\raggedright\arraybackslash}X@{}}
\toprule
\textbf{ID} & \textbf{Condition} & \textbf{Why it exists} \\
\midrule

\cref{res:monotone}
 & $t$ nondecreasing along $\po$ within each process
 & The clock is a lower bound on all future origin times. Both promotion
   guards conclude ``$p$ can no longer produce this'' from ``$p$'s clock has
   passed it''. \\

\cref{res:lookahead}
 & put: $\delta \ge \lag{p}{q}$
 & Makes \textsc{Publish} sound: a neighbour known to be at $\kn{p}{q}$ cannot
   write earlier than $\kn{p}{q} + \lag{p}{q}$. \\

\cref{res:staleness}
 & get: $|\delta| \le \lag{p}{q}$
 & Makes \textsc{Retire} sound: a neighbour cannot read further back than
   $\kn{p}{q} - \lag{p}{q}$. \\

\cref{res:phase}
 & within a process, every put on $(q,\eff)$ precedes every get on $(q,\eff)$
 & Closes the single gap left by the non-strict clock clause $\clk{p} \ge
   \eff$: a local zero-delta write at the current tick. \\

\cref{res:write-disjoint}
 & no two puts share a cell
 & Makes $\StoreH$ single-valued. Without it the semantics does not exist,
   and results vary between runs. \\

\cref{res:horizon}
 & $\delta \le \delta_{\max}$ for puts
 & Bounds live versions. Affects no value, only memory. \\

\cref{res:collective}
 & collectives called by all processes in the same order
 & Session, barrier and structure creation are collective. \\

\cref{res:terminate}
 & $\ell \ge \max_o \opeff{o}$ globally, identical on all processes
 & The barrier must outlast the furthest effective time in play. \\

\bottomrule
\end{tabularx}
\caption{The eight restrictions.}
\label{tab:restrictions}
\end{table}

\section{Enforcement status and proposed checks}

None of the eight is enforced today. The table below proposes a check for each
and estimates its cost, on the assumption that they would be compiled in under
a debug configuration and out of production builds.

\begin{table}[htbp]
\centering
\small
\begin{tabularx}{\textwidth}{@{}l >{\raggedright\arraybackslash}p{2.5cm}
                                >{\raggedright\arraybackslash}X
                                >{\raggedright\arraybackslash}p{1.9cm}@{}}
\toprule
\textbf{ID} & \textbf{Now} & \textbf{Proposed check} & \textbf{Cost} \\
\midrule
\cref{res:monotone}       & none & reject $t < \clk{p}$ in put/get
  & $O(1)$ \\
\cref{res:lookahead}      & none & reject $\delta < \lag{p}{q}$; the matrix
  lookup is already performed nearby & $O(1)$ \\
\cref{res:staleness}      & none & reject $|\delta| > \lag{p}{q}$
  & $O(1)$ \\
\cref{res:phase}          & partial: \code{abort()} on apply to a
  non-writable version, which detects it late and only sometimes
  & keep, and document that the abort indicates this restriction
  & --- \\
\cref{res:write-disjoint} & none & per-version interval set of written byte
  ranges; abort on overlap & $O(\log w)$ per put \\
\cref{res:horizon}        & none & optional session-level cap on $\delta$
  & $O(1)$ \\
\cref{res:collective}     & partial: dgraph uses
  \code{dstm\_msg\_consensus}, darray does not & extend the consensus pattern
  to darray & one all-reduce \\
\cref{res:terminate}      & none & all-reduce $\max \eff$ inside the barrier
  and compare against $\ell$ & one all-reduce \\
\bottomrule
\end{tabularx}
\caption{Proposed enforcement. Only \cref{res:write-disjoint} costs more than
a constant per operation, and it is the one whose violation is otherwise
undetectable.}
\label{tab:enforcement}
\end{table}

\section{Which results depend on which restriction}

\begin{table}[htbp]
\centering
\small
\begin{tabularx}{\textwidth}{@{}>{\raggedright\arraybackslash}p{5.2cm}
                                >{\raggedright\arraybackslash}X@{}}
\toprule
\textbf{Result} & \textbf{Consumes} \\
\midrule
\cref{lem:store-welldef} (store is well defined)
 & \cref{res:write-disjoint} \\
\cref{thm:write-closure} (write closure)
 & \cref{res:monotone}, \cref{res:lookahead}, \cref{res:phase}; (G2), (G3) \\
\cref{thm:immutability} (immutability)
 & via \cref{thm:write-closure} \\
\cref{thm:discard-safety} (discard safety)
 & \cref{res:monotone}, \cref{res:staleness}; (G2) \\
\cref{thm:read-soundness} (reads match $\StoreH$)
 & \cref{thm:write-closure}, \cref{res:write-disjoint} \\
\cref{thm:determinism} (determinacy)
 & all of the mechanical group, plus \cref{res:write-disjoint} \\
\cref{thm:deadlock-freedom} (no deadlock)
 \& (G3); fairness \ref{f:delivery}--\ref{f:clock} \\
\cref{thm:barrier-correct} (termination)
 & \cref{res:collective}, \cref{res:terminate} \\
\cref{thm:acyclic} (schedulability)
 & \cref{res:lookahead}, \cref{res:phase}; (G3) \\
\cref{thm:space-bound} (memory bound)
 & \cref{res:horizon}, \cref{thm:discard-safety} \\
\bottomrule
\end{tabularx}
\caption{Dependency of the results on the restrictions. (G3) --- every
channel lag is at least one --- appears in three separate results and
is the single most load-bearing well-formedness condition.}
\label{tab:result-deps}
\end{table}
